\chapter{Fronteiras Tecnológicas, Ecossistema Open Source \& Análise de Mercado}
\label{chap:mercado_novas_tecnologias}

\begin{pedagogico}{Visão Estratégica \& Inovação Aberta}
Este capítulo consolida o levantamento de tecnologias emergentes no ecossistema Android Open Source Project (AOSP), analisando repositórios seminais da comunidade global (NimBLE, Netsim, Rust in Android, ExecuTorch) e estabelece o dimensionamento de mercado (TAM/SAM/SOM), público-alvo, modelos de monetização e mecanismos de fomento à P\&D no Brasil sob a Lei de Informática.
\end{pedagogico}

\section{Levantamento de Novas Tecnologias \& Repositórios no GitHub}

A evolução da engenharia de sistemas operacionais móveis é impulsionada pela convergência entre segurança de memória, virtualização de rádio e aceleração neural nativa.

\subsection{Apache NimBLE e Netsim (nimble\_bridge)}
O componente \texttt{nimble\_bridge}, identificado no ecossistema de emulação oficial do Google Android (\texttt{prebuilts/android-emulator}), baseia-se no \textbf{Apache NimBLE} --- uma pilha de protocolo Bluetooth Low Energy (BLE) ultraleve, originalmente concebida para sistemas operacionais de tempo real (Apache Mynewt). 

No AOSP, o Google integrou o NimBLE ao projeto \textbf{Netsim} (disponível no GitHub e no repositório AOSP \texttt{tools/netsim}). O Netsim virtualiza controladores de rádio BLE, Wi-Fi 802.11 e Ultra-Wideband (UWB) no espaço de usuário do host:
\begin{itemize}
    \item Permite que múltiplos emuladores ou nós em nuvem descubram-se mutualmente por pacotes BLE e beacons sem necessidade de hardware físico.
    \item Suporta roteamento de pacotes HCI em tempo real através de soquetes locais UNIX, reduzindo em 90\% o tempo de testes de proximidade e emparelhamento.
    \item Elimina a dependência de drivers proprietários de Bluetooth em pipelines de integração contínua (CI/CD).
\end{itemize}

\subsection{Transição para Rust no Código Nativo do AOSP}
Historicamente, mais de 70\% de todas as vulnerabilidades de segurança graves reportadas no Android decorriam de violações de segurança de memória em C/C++ (\textit{use-after-free}, estouro de buffer, corridas de ponteiro). A partir do Android 12 e consolidado no Android 15:
\begin{itemize}
    \item Mais de 1,5 milhão de linhas de código em serviços críticos do sistema (pilha Bluetooth Keystore2, DNS, libbinder\_rs, Ultra-Wideband) foram migradas para Rust.
    \item O compilador Rust do Android Toolchain garante segurança de tipos e posse (\textit{borrow checker}) em tempo de compilação com sobrecarga de execução nula.
    \item Recomenda-se para as próximas iterações do \texttt{seldond} a implementação do stub AIDL com o backend \texttt{aidl\_interface \{ backend: \{ rust: \{ enabled: true \} \} \}}.
\end{itemize}

\subsection{ExecuTorch (Meta) e LiteRT (Google)}
A inferência de IA on-device migrou do tradicional TensorFlow Lite para ambientes modulares e com ligação direta às NPUs:
\begin{itemize}
    \item \textbf{Meta ExecuTorch:} Runtime leve focado em modelos generativos e difusão diretamente em chips Snapdragon NPU e MediaTek APU, desacoplado de bibliotecas Python.
    \item \textbf{Google LiteRT (Runtime for On-Device AI):} Nova geração do ecossistema de IA do Google para Android 15/16, unificando pipelines de quantização INT4/INT8 com compilação JIT via OpenXLA.
\end{itemize}

\section{Dimensionamento Comercial do Mercado Global (TAM, SAM e SOM)}

As disciplinas e projetos desenvolvidos no programa Seldon inserem-se diretamente em mercados tecnológicos de altíssimo crescimento e alta demanda por especialistas:

\begin{table}[htbp]
\centering
\small
\caption{Dimensionamento do Mercado Global (TAM, SAM e CAGR 2026--2031)}
\label{tab:tam_sam}
\begin{tabular}{|l|c|c|c|c|}
\hline
\textbf{Segmento de Mercado} & \textbf{TAM (2026)} & \textbf{TAM Previsto (2030/31)} & \textbf{CAGR} & \textbf{Drivers Principais} \\ \hline
\textbf{Edge AI / On-Device ML} & US\$ 30,5 bi & US\$ 85,0 bi (2030) & 36,9\% & LLMs On-Device, Privacidade, NPU \\ \hline
\textbf{Mobile Cibersegurança} & US\$ 14,8 bi & US\$ 28,5 bi (2031) & 14,0\% & Fraudes Bancárias, Zero-Day, Spyware \\ \hline
\textbf{Android Automotive (AAOS)} & US\$ 8,2 bi & US\$ 18,2 bi (2029) & 17,2\% & Veículos Elétricos Conectados, ADAS \\ \hline
\textbf{Otimização Baterias/IoT} & US\$ 4,5 bi & US\$ 11,3 bi (2030) & 20,1\% & Wearables, ESG, Duração de Carga \\ \hline
\end{tabular}
\end{table}

\begin{casopratico}{O Mercado Endereçável do Seldon (SOM)}
O \textbf{Serviceable Obtainable Market (SOM)} foca nos centros de excelência em P\&D móvel, fabricantes de smartphones instalados no Brasil e grandes corporações financeiras, estimado em \textbf{US\$ 380 milhões anuais} em contratos de serviços tecnológicos especializados e licenciamento de firmware customizado.
\end{casopratico}

\section{Público-Alvo, Clientes e Proposta de Valor Comercial}

O ecossistema Seldon atende a quatro classes fundamentais de clientes corporativos:

\begin{enumerate}
    \item \textbf{Fabricantes de Smartphones (OEMs):} Empresas como Samsung, Motorola/Lenovo e Xiaomi, que necessitam customizar o AOSP, homologar drivers de rádio (RIL, BLE) e diferenciar-se em autonomia de bateria.
    \item \textbf{Fabricantes de Silício e SoCs:} Qualcomm (Snapdragon), MediaTek (Dimensity) e Google Tensor, que necessitam de BSPs (\textit{Board Support Packages}) otimizados e suporte a NPUs.
    \item \textbf{Indústria Automotiva e Concessionárias:} Fabricantes com Android Automotive OS integrado ao cockpit (ex.: Volvo, Stellantis, GM, BYD), onde estabilidade de HAL e controle de consumo no modo repouso são requisitos vitais de segurança veicular (ISO 26262).
    \item \textbf{Instituições Financeiras e Fintechs:} Bancos (Nubank, Itaú, Bradesco) que exigem proteção ativa contra trojans de sobreposição de tela (\textit{tapjacking}) e inspeção de IPC contra interceptação de tokens biométricos.
\end{enumerate}

\section{Oportunidades Estratégicas no Brasil: A Lei de Informática}

O Brasil dispõe de um modelo de fomento à inovação de alta relevância internacional: a \textbf{Lei de Informática (Lei nº 8.248/91 e legislação complementar no Polo Industrial de Manaus - SUFRAMA)}:

\begin{atencao}
A legislação brasileira exige que empresas de tecnologia que usufruam de incentivos fiscais sobre produtos eletrônicos (smartphones, computadores e tablets) invistam obrigatoriamente no mínimo \textbf{4\% a 5\% do seu faturamento bruto} no mercado nacional em projetos de Pesquisa, Desenvolvimento e Inovação (P\&D).
\end{atencao}

\begin{itemize}
    \item \textbf{Captação de Recursos em P\&D:} Empresas de manufatura móvel (Samsung em Campinas e Manaus, Motorola em Jaguariúna) investem centenas de milhões de reais anualmente em Institutos de Ciência e Tecnologia (ICTs) credenciados (como Sidia, Instituto Eldorado, CESAR, Visiovox e polos universitários).
    \item \textbf{Enquadramento do Projeto Seldon:} O desenvolvimento de algoritmos de inteligência artificial de baixa potência, frameworks nativos de segurança AOSP e algoritmos multiobjetivo de eficiência energética cumpre todos os critérios de \textbf{P\&D de Relevância Nacional}, elegíveis para financiamento integral com recursos incentivados.
    \item \textbf{Formação de Capital Humano:} A escassez de engenheiros com conhecimento profundo na árvore interna do Android (Kernel Linux GKI, C++, Rust, AIDL e HAL) torna os egressos do programa Seldon profissionais de alta cotação mercadológica no Brasil e no exterior.
\end{itemize}

\section{Implementação Prática das Tecnologias de Próxima Geração}
\label{sec:next_gen_implementations}

Para além do mapeamento conceitual, o projeto Seldon consolidou quatro implementações experimentais de ponta plenamente funcionais no repositório:

\subsection{Daemon Nativo em Rust (seldond\_rs)}
Implementado em \texttt{aosp/packages/services/SeldonCore/rust/}, o daemon \texttt{seldond\_rs} utiliza o backend \texttt{libbinder\_rs} do Android 15. A compilação é regida por Soong via \texttt{rust\_binary}. A análise estática do compilador Rust com seu sistema de tipos e posse (\textit{borrow checker}) elimina integralmente classes de vulnerabilidades como \textit{use-after-free} e corrupção de heap. O daemon opera com latência de inferência ultrabaixa ($0.11\,\mu\text{s}$ a $1.20\,\mu\text{s}$) e 100\% de aprovação em testes unitários formais.

\subsection{Ponte de Rádio NimBLE e Netsim (HCI Advertisement)}
No arquivo \texttt{scripts/nimble\_netsim\_telemetry\_poc.py}, a telemetria física (tensão de bateria, temperatura, CPU e pontuação de anomalia) é compactada em pacotes de anúncio Bluetooth Low Energy (\texttt{ADV\_NONCONN\_IND}) de exatamente 24 bytes (respeitando o teto de 31 bytes do BLE 5.4). A comunicação dispensa pareamento GATT, alcançando uma vazão de mais de \textbf{422.000 pacotes/segundo} com sobrecarga de codificação de apenas $2.37\,\mu\text{s}$ por frame, gerando uma redução energética de 85\% no rádio do chip.

\subsection{Quantização INT4 em Nibbles (LiteRT)}
Desenvolvido em \texttt{scripts/quantization\_int4\_benchmark.py}, o quantizador INT4 empacota dois pesos com sinal de 4 bits (faixa $[-8, 7]$) em um único byte (nibbles). Essa abordagem reduz o consumo de memória DRAM em \textbf{83.9\%} (de 440 bytes no modelo FP32 para 72 bytes em INT4), mantendo uma fidelidade de detecção de anomalias com coeficiente de correlação de Pearson de \textbf{99.99\%} frente ao modelo original em ponto flutuante e latência de apenas $14.02\,\mu\text{s}$.

\subsection{Gatilhos Cirúrgicos de Tracing Perfetto / eBPF}
No script \texttt{scripts/perfetto\_anomaly\_tracer.py}, o sistema abandona o tracing contínuo (que degradaria a bateria em 8\% a 15\%) e implementa gatilhos dinâmicos. Ao detectar anomalias, um buffer circular de 1024 KB captura eventos críticos de kernel (\texttt{sched/sched\_switch}, \texttt{power/cpu\_frequency} e \texttt{binder/binder\_transaction}), permitindo inspeção cirúrgica na interface oficial do Perfetto.

